\section{总体架构设计}

\subsection{总体架构说明}
系统总体上采用“前后端分离 + 统一业务服务 + 异步任务调度 + 独立AI推理服务”的组合架构。该架构不是脱离现有工程基础的重新假设，而是在当前项目代码已经形成用户端、管理端、Django 后端、Celery 异步任务、Redis 消息中转和独立图像推理服务的基础上，对未来迭代目标进行整理和提升后的结果。

从业务演进方向看，系统已经不再只是处理图像伪造检测，而是需要统一承载学术图像、全篇论文和 Review 文本三类检测对象。因此，体系结构必须从“围绕图像功能组织实现”过渡为“围绕统一资源、统一任务和统一结果组织实现”。前端负责交互与展示，后端负责业务编排，异步任务系统负责长耗时检测任务的排队和执行，AI 推理服务负责具体模型运行，数据库与文件存储负责持久化和追溯。通过这种职责拆分，可以在保留现有图像链路的同时，为论文检测、Review 检测和综合报告生成预留稳定扩展点。

需要说明的是，本章为了强调系统静态结构和依赖关系，统一采用 UML 类图风格描述主要子系统和核心抽象，而不使用传统流程框图。这里的“类”并不局限于面向对象代码类，而是作为对系统中稳定职责单元的结构化表达。

\subsection{架构选型依据}
本系统选择当前总体架构，主要基于以下考虑：
\begin{itemize}[leftmargin=2em]
    \item \textbf{Web 访问是主要交互方式}。系统主要面向浏览器端用户，用户端和管理端存在明显不同的页面需求与权限边界，采用前后端分离结构更有利于界面独立演进；
    \item \textbf{检测任务执行时间较长}。无论是图像检测、论文检测还是 Review 检测，都可能涉及文件解析、特征提取、批量切分和模型推理，不适合同步阻塞式请求处理；
    \item \textbf{检测对象存在明显差异}。图像、论文和 Review 的输入格式、预处理方式、结果结构和模型调用方式并不一致，需要通过统一抽象之上的适配机制来隔离差异；
    \item \textbf{系统需要持续迭代扩展}。本轮只是综合学术 AI 鉴伪平台的第一阶段扩展，后续仍可能增加模型管理、报告策略、组织治理和更细粒度的人工复核能力，因此架构必须支持增量演进；
    \item \textbf{当前代码已具备可复用基础}。继续整理和提升现有 Django、Celery、Redis、WebSocket 与独立 AI 服务体系，相比完全推翻重写更具工程可行性。
\end{itemize}

\subsection{系统总体静态结构类图}
\begin{figure}[H]
    \centering
    \resizebox{0.98\textwidth}{!}{\begin{tikzpicture}[
    node distance=2.0cm and 1.9cm,
    every node/.style={font=\small},
    umlclass/.style={
        draw,
        rectangle split,
        rectangle split parts=3,
        rectangle split part align={center,left,left},
        rounded corners=3pt,
        line width=0.8pt,
        minimum width=4.25cm,
        text width=4.25cm,
        inner sep=5pt
    },
    relation/.style={-{Latex[length=2.4mm]}, thick, draw=black!75}
]

\node[umlclass, fill=blue!8] (user) {
\textbf{用户端}
\nodepart{second}
+ 资源上传\\
+ 检测发起\\
+ 结果查看
\nodepart{third}
+ 提交资源()\\
+ 创建任务()\\
+ 查询结果()
};

\node[umlclass, fill=blue!8, right=of user] (admin) {
\textbf{管理端}
\nodepart{second}
+ 用户治理\\
+ 审核审批\\
+ 日志统计
\nodepart{third}
+ 审批复核()\\
+ 管理组织()\\
+ 查询日志()
};

\node[umlclass, fill=orange!12, below=of $(user)!0.5!(admin)$] (backend) {
\textbf{后端业务服务}
\nodepart{second}
+ 认证鉴权\\
+ 任务编排\\
+ 报告生成
\nodepart{third}
+ 身份校验()\\
+ 分派任务()\\
+ 生成报告()
};

\node[umlclass, fill=orange!10, below left=of backend] (task) {
\textbf{异步任务调度}
\nodepart{second}
+ 任务排队\\
+ 状态同步\\
+ 失败重试
\nodepart{third}
+ 加入队列()\\
+ 执行任务()\\
+ 更新状态()
};

\node[umlclass, fill=orange!10, below right=of backend] (ai) {
\textbf{AI推理服务}
\nodepart{second}
+ 图像检测\\
+ 论文检测\\
+ 评审检测
\nodepart{third}
+ 检测图像()\\
+ 检测论文()\\
+ 检测评审()
};

\node[umlclass, fill=green!10, below=of task] (db) {
\textbf{数据库}
\nodepart{second}
+ 用户数据\\
+ 任务数据\\
+ 审核日志
\nodepart{third}
+ 保存()\\
+ 查询()\\
+ 更新()
};

\node[umlclass, fill=green!10, below=of ai] (fs) {
\textbf{文件存储}
\nodepart{second}
+ 上传文件\\
+ 中间产物\\
+ 报告文件
\nodepart{third}
+ 存储()\\
+ 读取()\\
+ 归档()
};

\draw[relation] (user) -- node[left, font=\scriptsize, fill=white, inner sep=1pt] {接口访问} (backend);
\draw[relation] (admin) -- node[right, font=\scriptsize, fill=white, inner sep=1pt] {接口访问} (backend);
\draw[relation] (backend) -- node[left, font=\scriptsize, fill=white, inner sep=1pt] {任务分派} (task);
\draw[relation] (task) -- node[above, font=\scriptsize, fill=white, inner sep=1pt] {调用} (ai);
\draw[relation] (backend) -- node[left, font=\scriptsize, fill=white, inner sep=1pt] {读写} (db);
\draw[relation] (backend) -- node[right, font=\scriptsize, fill=white, inner sep=1pt] {读写} (fs);
\draw[relation] (task) -- node[left, font=\scriptsize, fill=white, inner sep=1pt] {状态回写} (db);
\draw[relation] (ai) -- node[right, font=\scriptsize, fill=white, inner sep=1pt] {结果产物} (fs);
\draw[relation] (task) -- node[below, font=\scriptsize, fill=white, inner sep=1pt] {回调} (backend);

\end{tikzpicture}
}
    \caption{系统总体静态结构 UML 类图}
\end{figure}

图 3-1 从系统静态依赖关系角度展示了本项目的总体结构。用户端和管理端分别面向普通用户与管理用户，通过统一后端业务服务访问平台能力；后端业务服务一方面直接管理认证鉴权、组织与资源治理、结果查询和报告生成等同步能力，另一方面通过异步任务调度子系统将长耗时检测请求提交到 AI 推理服务；推理结果再由业务服务统一汇总并落入数据库与文件存储，最终由通知与状态反馈机制返回到前端界面。

与传统的分层框图相比，使用类图样式表达的好处在于可以更清楚地表示各个稳定子系统之间的静态关联关系、依赖方向和职责边界，便于后续详细设计时继续向模块类图和实体类图细化。

\subsection{统一资源与检测扩展关系类图}
\begin{figure}[H]
    \centering
    \resizebox{0.98\textwidth}{!}{\begin{tikzpicture}[
    node distance=1.9cm and 1.8cm,
    every node/.style={font=\small},
    umlclass/.style={
        draw,
        rectangle split,
        rectangle split parts=3,
        rectangle split part align={center,left,left},
        rounded corners=3pt,
        line width=0.8pt,
        minimum width=4.15cm,
        text width=4.15cm,
        inner sep=5pt
    },
    relation/.style={-{Latex[length=2.4mm]}, thick, draw=black!75},
    inherit/.style={-{Triangle[open,length=3.2mm,width=4.8mm]}, thick, draw=black!75}
]

\node[umlclass, fill=blue!10] (resource) {
\textbf{学术资源}
\nodepart{second}
+ 资源编号\\
+ 资源类型\\
+ 所有者编号
\nodepart{third}
+ 创建()\\
+ 校验()\\
+ 归档()
};

\node[umlclass, fill=blue!10, right=of resource] (task) {
\textbf{检测任务}
\nodepart{second}
+ 任务编号\\
+ 任务状态\\
+ 创建时间
\nodepart{third}
+ 提交()\\
+ 调度()\\
+ 完成()
};

\node[umlclass, fill=blue!10, right=of task] (result) {
\textbf{检测结果}
\nodepart{second}
+ 结果编号\\
+ 风险等级\\
+ 报告引用
\nodepart{third}
+ 汇总()\\
+ 持久化()\\
+ 导出报告()
};

\node[umlclass, fill=orange!10, below=of task] (adapter) {
\textbf{<<抽象>> 检测适配器}
\nodepart{second}
+ 资源类型\\
+ 模型配置
\nodepart{third}
+ 预处理()\\
+ 推理()\\
+ 后处理()
};

\node[umlclass, fill=orange!8, below left=of adapter] (image) {
\textbf{图像检测适配器}
\nodepart{second}
+ 图像元数据\\
+ 可疑区域
\nodepart{third}
+ 检测伪造()\\
+ 生成热力图()
};

\node[umlclass, fill=orange!8, below=of adapter] (paper) {
\textbf{论文检测适配器}
\nodepart{second}
+ 分段结果\\
+ AIGC评分
\nodepart{third}
+ 切分论文()\\
+ 检测AIGC()
};

\node[umlclass, fill=orange!8, below right=of adapter] (review) {
\textbf{评审检测适配器}
\nodepart{second}
+ 评审文本\\
+ 模板评分
\nodepart{third}
+ 检测模板化()\\
+ 检测生成痕迹()
};

\node[umlclass, fill=green!10, below=of result] (report) {
\textbf{报告服务}
\nodepart{second}
+ 报告模板\\
+ 导出格式
\nodepart{third}
+ 生成报告()\\
+ 合并结果()
};

\draw[relation] (resource) -- node[above, font=\scriptsize, fill=white, inner sep=1pt] {创建} (task);
\draw[relation] (task) -- node[above, font=\scriptsize, fill=white, inner sep=1pt] {产出} (result);
\draw[relation] (task) -- node[right, font=\scriptsize, fill=white, inner sep=1pt] {使用} (adapter);
\draw[inherit] (image) -- (adapter);
\draw[inherit] (paper) -- (adapter);
\draw[inherit] (review) -- (adapter);
\draw[relation] (result) -- node[right, font=\scriptsize, fill=white, inner sep=1pt] {导出} (report);

\end{tikzpicture}
}
    \caption{统一资源与检测扩展关系 UML 类图}
\end{figure}

图 3-2 展示了本轮迭代在架构抽象层面最关键的一次调整，即将现有系统中偏图像中心的实现方式，提升为统一资源与统一检测任务的结构。具体而言，系统通过 \textbf{AcademicResource} 抽象统一承载图像、论文和 Review 三类对象，通过 \textbf{DetectionTask} 抽象统一管理检测请求和状态流转，通过 \textbf{DetectionResult} 抽象统一组织检测输出，并通过 \textbf{DetectionAdapter} 抽象隔离不同资源类型对应的模型调用差异。

这一设计的直接收益在于：后续新增一种检测对象时，不必改动整个任务链路和结果链路，只需要在统一任务框架下补充新的资源表示和适配器实现即可。这种方式比在现有图像任务逻辑上叠加条件分支更稳定，也更适合课程项目后续继续迭代。

\subsection{系统组成}
从逻辑结构上看，系统由以下几个核心组成部分构成：
\begin{itemize}[leftmargin=2em]
    \item \textbf{用户端}：面向普通用户和审核参与者，负责资源上传、检测任务发起、结果查看、人工审核参与和通知展示；
    \item \textbf{管理端}：面向组织管理员和超级管理员，负责用户治理、组织管理、审核审批、日志查看、统计分析和模型配置入口；
    \item \textbf{后端业务服务}：负责认证鉴权、资源管理、检测任务创建、结果聚合、报告生成、权限控制和管理接口暴露；
    \item \textbf{异步任务调度子系统}：负责长耗时任务的排队、拆分、执行、失败重试和状态同步；
    \item \textbf{AI 推理子系统}：负责图像、论文和 Review 等不同检测模型的具体执行；
    \item \textbf{消息与通知子系统}：负责检测状态推送、系统通知分发和用户消息反馈；
    \item \textbf{数据库与文件存储}：负责业务数据、日志、上传资源、检测中间产物和报告文件的持久化。
\end{itemize}

\subsection{子系统职责划分}
\begin{longtable}{|p{3.5cm}|p{9.3cm}|p{2.2cm}|}
\hline
子系统 & 主要职责 & 典型技术实现 \\ \hline
用户端子系统 & 面向普通用户和审核用户提供资源上传、检测发起、结果展示、审核参与和通知接收能力 & Vue + Vuetify \\ \hline
管理端子系统 & 面向组织管理员和超级管理员提供治理、统计、日志查询和审核审批能力 & Vue + Vuetify \\ \hline
业务服务子系统 & 提供统一 REST 接口，承担认证鉴权、权限控制、业务编排、结果封装和报告管理 & Django + DRF \\ \hline
消息与实时通知子系统 & 负责 WebSocket 推送、任务状态广播、系统公告与用户通知分发 & Django Channels \\ \hline
异步调度子系统 & 负责长时任务调度、批处理拆分、失败重试、执行收尾和状态回写 & Celery + Redis \\ \hline
AI 推理子系统 & 负责检测模型推理、特征计算、风险分析和差异化结果生成 & 独立 Python 推理服务 \\ \hline
数据存储子系统 & 负责业务数据、任务状态、日志记录、审核记录和元数据存储 & MySQL \\ \hline
文件存储子系统 & 负责上传资源、结果文件、可视化产物和报告文件保存 & 本地媒体目录/文件系统 \\ \hline
\end{longtable}

\subsection{核心调用链设计}
系统的核心调用链可概括为以下几个阶段：
\begin{enumerate}[leftmargin=2em]
    \item 用户通过用户端或管理端发起业务请求，例如上传资源、提交检测、申请人工复核或审批审核任务；
    \item 前端通过 HTTP API 或 WebSocket 将请求发送到后端业务服务；
    \item 后端完成身份认证、权限校验、资源登记、任务创建和状态初始化，并根据任务性质决定同步处理还是异步处理；
    \item 对于检测类长时请求，后端将任务提交到异步任务调度子系统，由调度子系统负责排队、分发和失败重试；
    \item 异步任务调度子系统根据任务类型调用对应的检测适配器，由适配器进一步调用 AI 推理服务完成模型执行；
    \item 推理结果经过后端业务服务统一汇总后写入数据库与文件存储，并触发报告生成、通知创建和状态变更；
    \item 前端通过轮询接口或实时推送机制获取任务最新状态、检测结果和报告链接。
\end{enumerate}

上述调用链将“用户请求响应速度”和“检测处理耗时”有效解耦，使平台既能保证交互上的及时反馈，又能处理较重的文件读写、模型推理和报告生成任务。

\subsection{分层设计}
为了使系统在结构上更清晰、在实现上更易维护，本系统采用分层设计。各层职责如下表所示：
\begin{longtable}{|p{3cm}|p{11cm}|}
\hline
层次 & 职责 \\ \hline
表示层 & 提供用户端与管理端页面、表单交互、结果展示和状态反馈，仅负责界面组织与前端状态管理，不直接承载核心业务规则。\\ \hline
接口层 & 提供 HTTP API 与 WebSocket 接口，对外暴露统一访问入口，负责请求接收、参数校验和基础异常响应。\\ \hline
业务层 & 实现用户、组织、资源、检测、审核、报告、通知和统计等核心业务逻辑，是系统的业务中枢。\\ \hline
任务层 & 处理异步检测任务、批处理调度、状态同步、任务收尾和报告生成触发。\\ \hline
适配层 & 对不同资源类型对应的检测模型和外部推理能力进行适配与封装，屏蔽底层模型差异。\\ \hline
数据层 & 存储用户数据、任务数据、审核数据、日志、文件路径和报告元数据，保证业务数据可持久化和可追溯。\\ \hline
\end{longtable}

\subsection{架构约束与关键设计决策}
为了保证后续详细设计和编码实现能够沿同一方向演进，本系统在总体架构层面明确以下约束与设计决策：
\begin{itemize}[leftmargin=2em]
    \item \textbf{前端不直接调用模型服务}。所有检测请求必须先进入后端业务服务，由后端统一完成权限校验、任务登记和状态管理；
    \item \textbf{检测必须通过统一任务中心流转}。无论检测对象是图像、论文还是 Review，原则上都应通过统一的任务创建、调度和回写机制执行；
    \item \textbf{结果主体与差异化载荷分离}。统一结果对象只保存通用状态、风险等级和报告引用，具体差异化内容由资源类型专属结构承载；
    \item \textbf{模型接入通过适配器隔离}。业务层不直接依赖具体模型实现，模型替换和模型扩展应限制在适配层及推理子系统内部；
    \item \textbf{人工审核链路独立于自动检测链路存在}。人工审核可以引用自动检测结果，但不应与模型执行过程耦合，从而保证人工复核流程可单独演进；
    \item \textbf{关键过程必须可追溯}。资源上传、任务创建、状态变化、审核结论和报告生成都需要在数据层保留可回查记录。
\end{itemize}

\subsection{面向迭代的架构调整方向}
考虑到当前实现仍以图像检测流程为中心，架构上需要逐步向统一资源架构过渡。该过渡主要体现在以下几个方面：
\begin{itemize}[leftmargin=2em]
    \item \textbf{资源对象统一化}：从仅支持图像的资源表示过渡为统一支持图像、论文和 Review 的资源抽象；
    \item \textbf{检测能力适配化}：图像检测、论文检测和 Review 检测分别通过独立适配器接入，避免在业务代码中硬编码不同检测逻辑；
    \item \textbf{结果结构层次化}：采用统一结果主体与差异化结果载荷相结合的方案，使不同资源类型既能纳入统一任务流程，又能保留各自结果特点；
    \item \textbf{治理能力泛化}：管理端从围绕图像任务治理扩展为面向统一学术资源、统一日志和统一模型配置的治理结构；
    \item \textbf{报告生成标准化}：将现有偏图像结果的报告输出能力扩展为可面向多资源类型的统一报告服务。
\end{itemize}

\subsection{关键架构优势}
本轮总体架构设计具备以下优势：
\begin{itemize}[leftmargin=2em]
    \item 能复用现有前后端、异步任务与图像检测链路，降低迭代开发成本；
    \item 能通过统一任务体系纳入论文和 Review 两类新资源，避免形成多个割裂子系统；
    \item 能将长时检测任务从用户请求链路中解耦，提高系统交互流畅性；
    \item 能将模型推理能力与业务编排逻辑分离，便于后续替换或新增模型；
    \item 能为后续详细设计中的数据库重构、接口扩展、实体建模和流程设计提供稳定基础。
\end{itemize}

\subsection{总体架构小结}
综上，本系统总体架构并不是对现有实现的简单罗列，而是对当前代码基础、迭代需求目标和后续演进方向进行综合平衡后的结果。其核心思想是：在保留当前图像检测平台工程基础的同时，通过统一资源、统一任务、统一结果和适配化模型接入，将系统演进为可持续扩展的综合学术 AI 鉴伪平台。
